% Part: computability
% Chapter: recursive-functions
% Section: composition

\documentclass[../../../include/open-logic-section]{subfiles}

\begin{document}

\olfileid{cmp}{rec}{com}
\olsection{অপেক্ষকের মিশ্রণ}

$f$ ও~$g$ স্বাভাবিক সংখ্যার দুটি এক-স্থানীয় অপেক্ষক হলে তাদের মেশানো
যায়: $h(x) = f(g(x))$। নতুন অপেক্ষক~$h(x)$-কে $f$ ও~$g$ অপেক্ষক
থেকে \emph{মিশ্রণের} মাধ্যমে সংজ্ঞায়িত করা হয়। একাধিক আর্গুমেন্টের
অপেক্ষকের ক্ষেত্রেও এটিকে সাধারণীকরণ করতে চাই।

এটি করার একটি পদ্ধতি হলো: ধরি, $f$ একটি $k$-স্থানীয় অপেক্ষক এবং
$g_0$, \dots, $g_{k-1}$ এমন $k$টি অপেক্ষক, যাদের প্রত্যেকটি
$n$-স্থানীয়। তাহলে নিচের মতো নতুন $n$-স্থানীয় অপেক্ষক~$h$
সংজ্ঞায়িত করতে পারি:
\[
h(x_0, \dots, x_{n-1}) =
f(g_0(x_0, \dots, x_{n-1}), \dots, g_{k-1}(x_0, \dots, x_{n-1}))
\]
$f$ ও সব $g_i$ গণনসাধ্য হলে~$h$-ও গণনসাধ্য। $h(x_0,
\dots, x_{n-1})$ গণনা করতে প্রথমে প্রতিটি $i=0$, \dots,~$k-1$-এর জন্য
$y_i = g_i(x_0, \dots, x_{n-1})$ মান গণনা করি। তারপর এই মানগুলি $f$-এ
দিয়ে $h(x_0, \dots, x_{k-1}) = f(y_0, \dots, y_{k-1})$ গণনা করি।

বিদ্যমান অপেক্ষক ব্যবহার করে নতুন অপেক্ষক গণনা করার সময় যা ঘটে, তার
এই বর্ণনা অতিরিক্ত সীমাবদ্ধ মনে হতে পারে। যেমন, কখনও কোনো অপেক্ষকের
সব আর্গুমেন্ট ব্যবহার করি না; আদিম পুনরাবৃত্তির মাধ্যমে~$\Add$
সংজ্ঞায়নের সময় $g(x, y, z) = \Succ(z)$ সংজ্ঞায়িত করার ক্ষেত্রটি
এমনই। নিচের অপেক্ষকগুলি ব্যবহার করার অনুমতি আছে ধরে নিই:
\[
\Proj{n}{i}(x_0, \dots, x_{n-1}) = x_i
\]
$\Proj{k}{i}$ অপেক্ষকগুলিকে \emph{অভিক্ষেপ অপেক্ষক} বলা হয়:
$\Proj{n}{i}$ হলো একটি $n$-স্থানীয় অপেক্ষক। তাহলে~$g$-কে এভাবে
সংজ্ঞায়িত করা যায়:
\[
g(x, y, z) = \Succ(\Proj{3}{2}(x, y, z)).
\]
এখানে $f$-এর ভূমিকা নিয়েছে $1$-স্থানীয় অপেক্ষক~$\Succ$, তাই $k=1$।
আর আমাদের একটি $3$-স্থানীয় অপেক্ষক~$\Proj{3}{2}$ আছে, যা $g_0$-এর
ভূমিকা পালন করে। ফলাফল হলো এমন একটি $3$-স্থানীয় অপেক্ষক, যা তৃতীয়
আর্গুমেন্টের উত্তরসূরি ফেরত দেয়।

অভিক্ষেপ অপেক্ষক আর্গুমেন্টগুলির বিন্যাস বদলে বা তাদের অভিন্ন করে নতুন
অপেক্ষক সংজ্ঞায়িত করতেও সাহায্য করে। যেমন, $h(x) = \Add(x, x)$ অপেক্ষকটি
এভাবে সংজ্ঞায়িত করা যায়:
\[
h(x_0) = \Add(\Proj{1}{0}(x_0),\Proj{1}{0}(x_0)).
\]
এখানে $k=2$, $n=1$; $f(y_0,y_1)$-এর ভূমিকা নিয়েছে~$\Add$, আর
$g_0(x_0)$ ও $g_1(x_0)$-এর ভূমিকা উভয়ই নিয়েছে
$\Proj{1}{0}(x_0)$, এক-স্থানীয় অভিক্ষেপ অপেক্ষকটি (অর্থাৎ অভিন্ন অপেক্ষক)।

আগে থেকে থাকা $f(y_0, y_1)$ অপেক্ষক দেওয়া থাকলে
$h(x_0, x_1) = f(x_1, x_0)$ অপেক্ষকটি সংজ্ঞায়িত করতে পারি:
\[
h(x_0, x_1) = f(\Proj{2}{1}(x_0, x_1),\Proj{2}{0}(x_0, x_1)).
\]
এখানে $k=2$, $n = 2$, এবং $g_0$ ও $g_1$-এর ভূমিকা যথাক্রমে
$\Proj{2}{1}$ ও~$\Proj{2}{0}$ পালন করে।

তোমার হয়তো আরও একটি আপত্তি হতে পারে যে $g_0$, \dots,~$g_{k-1}$-এর
সবার একই স্থানসংখ্যা~$n$ থাকা দরকার। (মনে রাখবে, কোনো অপেক্ষকের
\emph{স্থানসংখ্যা} হলো তার আর্গুমেন্টের সংখ্যা; $n$-স্থানীয় অপেক্ষকের
স্থানসংখ্যা~$n$।) কিন্তু অভিক্ষেপ অপেক্ষক যোগ করলে প্রয়োজনীয় নমনীয়তা
পাওয়া যায়। যেমন, ধরি $f$ ও~$g$ $3$-স্থানীয় অপেক্ষক এবং $h$ হলো নিচে
সংজ্ঞায়িত $2$-স্থানীয় অপেক্ষক:
\[
h(x,y) = f(x,g(x,x,y),y).
\]
$h$-এর সংজ্ঞাটি অভিক্ষেপ অপেক্ষক ব্যবহার করে এভাবে লেখা যায়:
\[
h(x,y) = f(\Proj{2}{0}(x,y), g(\Proj{2}{0}(x,y), \Proj{2}{0}(x,y),
\Proj{2}{1}(x,y)), \Proj{2}{1}(x,y)).
\]
তখন $h$ হলো $f$-এর সঙ্গে $\Proj{2}{0}$, $l$ এবং
$\Proj{2}{1}$-এর মিশ্রণ, যেখানে
\[
l(x,y) = g(\Proj{2}{0}(x,y),\Proj{2}{0}(x,y),\Proj{2}{1}(x,y)),
\]
অর্থাৎ $l$ হলো $g$-এর সঙ্গে $\Proj{2}{0}$, $\Proj{2}{0}$ এবং
~$\Proj{2}{1}$-এর মিশ্রণ।

\end{document}
